\begin{abstract}
If we read 15 documents from a corpus, how much of the corpus have we seen? Researchers who study text read samples because they cannot analyze every document. A survey sample comes with a formula for its size and a margin of error, and a sample of text comes with neither. We measure what a sample has missed as the percentage of the corpus with no document of the same content in the sample. A text encoder turns each document into an embedding, a list of numbers, so that documents with similar meaning have embeddings close together. Two documents count as the same content when their embeddings lie closer than a distance the researcher chooses, and two large language models judged pairs of documents to tell which distance means the same specific content and which means the same topic. We estimate the missed percentage from the sample alone with a classical estimator built on what we call the isolation rate, the percentage of sampled documents with no other sampled document of the same content. The main lesson is that a random sample is a poor use of a reading budget on corpora that repeat themselves, in which many documents make the same claim. With the whole corpus embedded, a researcher can divide it into as many groups of similar documents as the budget allows. The documents nearest the centres of these groups cover far more of such a corpus than the same number of random documents. At the level of topics, 15 documents chosen this way often miss about half as much of the corpus as 15 random documents. On corpora whose documents rarely share specific content, no choice of documents helps at the level of specific content, because a document unlike every other is covered only by reading it. The isolation rate of a random pilot of 100 documents tells a researcher in advance which of the two cases applies. To plan the size of a sample, we show that on most corpora we tested a larger sample and a looser definition of the same content trade against each other at a fixed rate. The rate is set by the intrinsic dimension of the corpus, the number of independent ways in which its documents differ in meaning. This rate gives a sample size formula. In 77 of 83 tests, each a combination of corpus, target and distance, the sample actually needed fell between two thirds and one and a half times the formula's answer. A larger pilot of 2{,}000 documents shows in advance whether the formula can be trusted for a given corpus.
\end{abstract}
